% ============================================================
% Abstract
% ============================================================

\begin{trabstractcn}
随着人工智能技术的快速发展，大语言模型（Large Language Models, LLMs）在多模态推理任务中展现出前所未有的潜力。本技术报告系统梳理了我们团队在基于大语言模型的多模态推理领域的最新研究进展。首先，我们提出了一种融合视觉-语言表征的渐进式推理框架（Progressive Multimodal Reasoning, PMR），通过分层注意力机制有效缓解模态间语义鸿沟问题。其次，我们在多个公开基准数据集上验证了所提方法的有效性，包括 ScienceQA、A-OKVQA 以及自建的 SIAT-MR 医学多模态推理数据集。实验结果表明，PMR 在复杂推理任务上的准确率较现有最优方法提升了 3.8\%。最后，我们讨论了当前方法的局限性，并展望了未来在具身智能和临床决策支持系统中的应用前景。
\trkeywordscn{大语言模型；多模态推理；视觉-语言融合；渐进式推理；临床决策支持}
\end{trabstractcn}

\vspace{1em}

\begin{trabstract}
With the rapid advancement of artificial intelligence, Large Language Models (LLMs) have demonstrated unprecedented potential in multimodal reasoning tasks. This technical report systematically summarizes our team's latest research progress in LLM-based multimodal reasoning. First, we propose a Progressive Multimodal Reasoning (PMR) framework that integrates vision-language representations via hierarchical attention mechanisms, effectively alleviating cross-modal semantic gaps. Second, we validate the proposed method on multiple public benchmarks, including ScienceQA, A-OKVQA, and our in-house SIAT-MR medical multimodal reasoning dataset. Experimental results show that PMR achieves a 3.8\% accuracy improvement over state-of-the-art methods on complex reasoning tasks. Finally, we discuss the limitations of current approaches and outline future applications in embodied intelligence and clinical decision support systems.
\trkeywords{Large Language Models; Multimodal Reasoning; Vision-Language Fusion; Progressive Reasoning; Clinical Decision Support}
\end{trabstract}

\newpage
